\ifdefined\isMainDocument
\else
    \documentclass{article}
    \usepackage{fullpage}
    % \usepackage{geometry}
    % \geometry{top=0.5in, bottom=0.5in, left=0.5in, right=0.5in}

    \usepackage{wrapfig}
    \usepackage{lineno}
    \usepackage{amsmath, amsthm, amsfonts, amssymb, amscd}
    \usepackage{bm, bbm, mathrsfs}
    \usepackage{mathtools}
    \usepackage{nicefrac}
    \usepackage{caption}
    % \usepackage{natbib}
    \usepackage{biblatex}
    \addbibresource{references.bib}
    \usepackage[colorinlistoftodos]{todonotes}
    \usepackage{graphicx}
    \graphicspath{figures/}
    \usepackage{setspace}
    \setstretch{1.5}

    \begin{document}
    \linenumbers
\fi
% Drafted 2026-10-23 (user) to the settled claim; under GENETICS's 250-word limit. Every number is
% from Results. Scaffolding; the voice is the author's to revise.

\section*{Abstract}
Pedigree-based animal models now estimate the additive genetic variance in lifetime fitness, $V_A$, in wild populations, where it is often substantial. Classical theory attaches two equilibrium expectations to such a variance: it should reduce neutral diversity, and it should be maintained by recurrent deleterious mutation at mutation--selection--drift balance. We put 19 pedigree estimates to both tests, under the multiplicative fitness model on which the estimates rest. We extend the reduction in effective size to social monogamy with extra-pair paternity and to linked sites along empirical linkage maps, and compute the variance that measured mutation rates and empirical distributions of fitness effects can maintain. The measured $V_A$ is too small to explain the neutral diversity these populations are missing and too large to be maintained by mutation. It roughly halves effective size in the median population, whereas the observed deficits of 230- to 38,000-fold would require a long-term $V_A$ above one. Maintaining it would require the average new mutation to lower heterozygous fitness by about 15\%, a median 18 times the mean effects estimated in related species. No rescaling of the estimates to a long-term value passes both tests: the required values differ about 390-fold on average and over 100-fold in every population. The conflict does not depend on the estimates, because a variance maintained by mutation would remove only a few percent of neutral diversity. Most of the variance measured in the wild must therefore be maintained by other processes, and most of the missing diversity has other causes.

\ifdefined\isMainDocument
\else
    \printbibliography
    \end{document}
\fi
